% Source PDF pages 127--130; manual pages 2-98--2-101.
\subsection{Search Methods}\label{sec:search-methods}

TAOS uses two classes of one-dimensional search. Root searches vary a parameter
$x$ until $f(x)=0$; minimization searches vary $x$ until $f(x)$ reaches a minimum.
Well-behaved one-dimensional searches may be nested to solve multidimensional
problems.

The selected method depends on the application. Segment final conditions use a
linear method because their search intervals are short enough that most functions are
nearly linear. Trajectory searches used to satisfy constraints generally use a
parabolic method because it often converges faster.

For root searches, the parameter is restricted to an interval
$x_{lo}<x<x_{hi}$. TAOS assumes that $f(x)$ is reasonably well behaved within the
interval and that a single root is present.

\taossubhead{Newton--Raphson Method}

The Newton--Raphson method starts with an estimate $x_n$ and approximates the
function locally by a straight line:
\begin{equation}
  f(x)=f(x_n)+(x-x_n)\left.\frac{\partial f}{\partial x}\right|_{x_n}.
  \label{eq:newton-linearization}
\end{equation}
Setting the linear approximation to zero gives the next estimate,
\begin{equation}
  x_{n+1}=x_n-
  \frac{f(x_n)}{\left.\dfrac{\partial f}{\partial x}\right|_{x_n}}.
  \label{eq:newton-update-search}
\end{equation}
The geometry is illustrated in \cref{fig:newton-raphson-search}. The derivative is
estimated numerically with a forward difference, and the update is repeated until
$f(x)$ is within the requested tolerance.

\begin{figure}[htbp]
  \centering
  \begin{tikzpicture}[>=Latex,x=1.0cm,y=0.9cm,line cap=round]
  \draw[->] (0,0)--(7.2,0) node[right] {$x$};
  \draw[->] (0,-1.2)--(0,4.1) node[above] {$f(x)$};
  \draw[line width=0.9pt]
    plot[smooth] coordinates {(0.45,3.5) (1.6,3.25) (2.7,2.65) (3.6,1.65) (4.25,0.45) (5.0,-0.55) (6.4,-1.05)};
  \coordinate (pn) at (3.55,1.72);
  \coordinate (root) at (4.45,0);
  \draw[densely dashed] (pn)--(3.55,0) node[below] {$x_n$};
  \draw[densely dashed] (pn)--(0,1.72) node[left] {$f(x_n)$};
  \draw[line width=0.75pt] (2.65,3.15)--(4.45,0)--(5.0,-0.95);
  \draw[->] (3.9,1.20)--(4.30,0.37);
  \node[right] at (3.9,1.45) {$\dfrac{\partial f}{\partial x}$};
  \fill (pn) circle (1.8pt);
  \fill (root) circle (1.8pt);
  \node[below] at (root) {$x_{n+1}$};
  \node[left] at (0,0) {$0$};
\end{tikzpicture}

  \caption{Newton--Raphson search.}
  \label{fig:newton-raphson-search}
\end{figure}

\taossubhead{Secant Method}

The secant method also assumes that $f(x)$ is linear over the current interval. Given
the function values at the lower and upper limits, the estimated root is
\begin{equation}
  x=x_{lo}-f(x_{lo})
  \left[
    \frac{x_{hi}-x_{lo}}{f(x_{hi})-f(x_{lo})}
  \right].
  \label{eq:secant-root-estimate}
\end{equation}

\begin{figure}[htbp]
  \centering
  \begin{tikzpicture}[>=Latex,x=1.0cm,y=0.9cm,line cap=round]
  \draw[->] (0,0)--(7.6,0);
  \draw (0,-1.15)--(0,3.55);
  \draw (7,-1.15)--(7,3.55);
  \draw[line width=0.9pt]
    plot[smooth] coordinates {(0,3.2) (1.35,2.9) (2.55,2.1) (3.45,0.45) (4.1,-0.75) (5.45,-1.15) (7,-1.45)};
  \draw[line width=0.75pt] (0,3.2)--(7,-1.45);
  \coordinate (xnew) at (4.82,0);
  \coordinate (fx) at (4.10,-0.75);
  \draw[densely dashed] (xnew)--(4.82,1.15);
  \draw[densely dashed] (fx)--(4.10,0);
  \node[align=center] at (4.72,1.45) {new\\$x_{hi}$};
  \draw[->] (5.25,2.05)--(6.65,2.05);
  \node[above,align=center] at (5.95,2.08) {Eliminate this part\\of the search interval};
  \node[below] at (0,-1.15) {$x_{lo}$};
  \node[below] at (7,-1.15) {$x_{hi}$};
  \node[below] at (4.10,-1.15) {$x$};
  \node[left] at (0,3.2) {$f(x_{lo})$};
  \node[right] at (6.75,-0.95) {$f(x_{hi})$};
  \node[left] at (fx) {$f(x)$};
  \node[left] at (0,0) {$0$};
\end{tikzpicture}

  \caption{Secant search.}
  \label{fig:secant-search}
\end{figure}

After evaluating $f(x)$, TAOS replaces the boundary whose function value has the
same sign as $f(x)$. Retaining endpoints of opposite sign preserves a bracket around
the root. The process continues until the estimated residual is within tolerance.

\taossubhead{Parabolic Root Search}

The parabolic method approximates the function by
\begin{equation}
  f(x)=a x^2+b x+c.
  \label{eq:parabolic-search-polynomial}
\end{equation}
Three function evaluations determine $a$, $b$, and $c$. During startup, TAOS
evaluates the function at $x-\Delta x$ and $x+\Delta x$, where $x$ is the initial
estimate and $\Delta x$ represents its expected accuracy. A linear estimate similar
to \cref{eq:secant-root-estimate} supplies the third point, with all points restricted
to the search interval. The geometry is shown in
\cref{fig:parabolic-root-search}.

\begin{figure}[htbp]
  \centering
  \begin{tikzpicture}[>=Latex,x=1.0cm,y=0.9cm,line cap=round]
  \draw[->] (0,0)--(8.0,0);
  \draw (0,-1.25)--(0,3.55);
  \draw (7.6,-1.25)--(7.6,3.55);
  \draw[line width=0.9pt]
    plot[smooth] coordinates {(0,3.1) (1.35,2.86) (2.7,2.15) (3.85,0.78) (4.85,-0.38) (6.1,-0.88) (7.6,-1.05)};
  \draw[line width=0.75pt,dashed]
    plot[smooth] coordinates {(1.0,2.98) (2.7,2.12) (3.95,0.82) (4.95,-0.05) (5.6,-0.28)};
  \coordinate (p1) at (1.25,2.88);
  \coordinate (p2) at (2.75,2.10);
  \coordinate (p3) at (4.95,-0.42);
  \draw[densely dashed] (p1)--(1.25,0);
  \draw[densely dashed] (p2)--(2.75,0);
  \draw[densely dashed] (p3)--(4.95,0);
  \fill (p1) circle (1.7pt);
  \fill (p2) circle (1.7pt);
  \fill (p3) circle (1.7pt);
  \node[below] at (1.25,0) {$x_1$};
  \node[below=0.34cm] at (1.25,0) {$x-\Delta x$};
  \node[below] at (2.75,0) {$x_2$};
  \node[below=0.34cm] at (2.75,0) {$x+\Delta x$};
  \node[below right] at (p3) {$x_3$};
  \node[align=left,text width=3.6cm] at (6.0,2.25)
    {A parabola is fit through the three points to obtain a new estimate of the root};
  \node[below] at (0,-1.25) {$x_{lo}$};
  \node[below] at (7.6,-1.25) {$x_{hi}$};
  \node[left] at (0,0) {$0$};
  \node[left] at (0,3.1) {$f(x)$};
\end{tikzpicture}

  \caption{Parabolic root search.}
  \label{fig:parabolic-root-search}
\end{figure}

For three points $x_1$, $x_2$, and $x_3$, the polynomial coefficients are
\begin{equation}
  b=
  \frac{
    (x_2^2-x_3^2)\,[f(x_1)-f(x_3)]
    -(x_1^2-x_3^2)\,[f(x_2)-f(x_3)]
  }{
    (x_2^2-x_3^2)(x_1-x_3)
    -(x_1^2-x_3^2)(x_2-x_3)
  },
  \label{eq:parabolic-search-b}
\end{equation}
\begin{equation}
  a=
  \frac{f(x_2)-f(x_3)-b(x_2-x_3)}{x_2^2-x_3^2},
  \label{eq:parabolic-search-a}
\end{equation}
\begin{equation}
  c=f(x_2)-a x_2^2-b x_2.
  \label{eq:parabolic-search-c}
\end{equation}
The two roots of the fitted parabola are
\begin{equation}
  x=\frac{-b\pm\sqrt{b^2-4ac}}{2a}.
  \label{eq:parabolic-search-roots}
\end{equation}

Normally only one root lies in the search interval. If both do, TAOS chooses the one
nearest the most recent best estimate; if that rule does not distinguish them, the
lower root is selected. After the new function evaluation, the point farthest from the
best root estimate is discarded, and the fit is repeated until the residual is within
tolerance.

\taossubhead{Golden-Section Search}

The golden-section method minimizes a function without derivatives. It evaluates
two interior points chosen to divide the current interval into golden sections:
\begin{equation}
  x_1=x_{hi}-\alpha(x_{hi}-x_{lo}),
  \label{eq:golden-section-x1}
\end{equation}
\begin{equation}
  x_2=x_{lo}+\alpha(x_{hi}-x_{lo}),
  \label{eq:golden-section-x2}
\end{equation}
where
$\alpha=(\sqrt{5}-1)/2$. Based on $f(x_1)$ and $f(x_2)$, one outer portion of the
interval is discarded, as illustrated in \cref{fig:golden-section-search}. One old
interior point can be reused after every interval reduction. The process stops when
$x_{hi}-x_{lo}$ is less than the requested tolerance.

\begin{figure}[htbp]
  \centering
  \begin{tikzpicture}[>=Latex,x=1.0cm,y=0.9cm,line cap=round]
  \draw[->] (0,0)--(7.8,0);
  \draw (0,0)--(0,3.15);
  \draw (7.2,0)--(7.2,3.15);
  \draw (2.25,0)--(2.25,1.05);
  \draw (5.05,0)--(5.05,2.55);
  \fill (2.25,1.05) circle (2pt);
  \fill (5.05,2.55) circle (2pt);
  \node[above left] at (2.25,1.05) {$f(x_1)$};
  \node[above left] at (5.05,2.55) {$f(x_2)$};
  \draw[<->] (0.25,2.18)--(4.80,2.18);
  \node[fill=white,align=center] at (2.55,2.18) {Minimum must be\\in this interval};
  \draw[->] (5.30,1.35)--(6.90,1.35);
  \node[below,align=center] at (6.10,1.32) {Eliminate this part\\of the search interval};
  \node[below] at (0,0) {$x_{lo}$};
  \node[below] at (2.25,0) {$x_1$};
  \node[below] at (5.05,0) {$x_2$};
  \node[below] at (7.2,0) {$x_{hi}$};
  \node[left] at (0,2.75) {$f(x)$};
\end{tikzpicture}

  \caption{Golden-section search.}
  \label{fig:golden-section-search}
\end{figure}

\taossubhead{Parabolic Minimization}

Parabolic minimization uses the same startup and coefficient formulas as the
parabolic root search, but evaluates the stationary point of the fitted parabola,
\begin{equation}
  x=-\frac{b}{2a}.
  \label{eq:parabolic-minimum}
\end{equation}
TAOS tests whether the stationary point is a minimum. If so, the interval is redefined
to retain that point near the middle; otherwise, the interval is redefined around the
three lowest function evaluations. The process repeats until the predicted value
$a x^2+b x+c$ agrees with the actual evaluation $f(x)$ within tolerance.
